\label{sec:negatives}

This section collects the results that did not come out the way we wanted.
They are part of the deliverable, not an embarrassment routed around.

\subsection{Retraction: the CNN Turing surrogate (Study A)}

An $n = 58$ pilot suggested a CNN could predict the beyond-onset wavenumber
upshift from the linear dispersion curve $17.7\%$ better than the linear
formula alone. At $n = 171$ the effect inverts (CNN rel.\ RMSE $0.292$--$0.295$
vs.\ linear $0.272$). The small-sample win was an artefact of the first
sample. What survives: the upshift phenomenon itself (Section~\ref{sec:upshift}),
the quantisation diagnosis, and a cautionary data point about
surrogate-modelling claims at $n < 100$. The interim claim was corrected in
the project record within a day; the retraction is committed in
\texttt{results/turing\_cnn.json} alongside the numbers.

\subsection{Negative: the second exponential component is not justified
per-embryo (Study C)}

Despite the population-level window sensitivity of $\lambda$
(Table~\ref{tab:windows}), per-embryo two-component fits reject the single
exponential in only $0.74\%$ of 676 gradients ($0.30\%$ Bonferroni) --- below
the 5\% false-positive floor. Whatever drives the window trend, it is too
small per embryo to license a two-component model. We report this because the
two-component fit is a standard move in this literature, and on this dataset
it fails an $F$-test at scale.

\subsection{Negative: uniform-sampling basin classification is degenerate
(Study B)}

A naive pipeline for ``learn the basin boundary'' samples states uniformly
and reports accuracy. Here that yields $99.3\%$ accuracy for the majority
classifier itself (Section~\ref{sec:basinlearn}): the number is meaningless
and would have looked publishable. Balanced sampling is what makes the
learning curve informative. Any future basin-learning benchmark on rare
attractors must control for this.

\subsection{Known limitation: the multi-cell segment-polarity extension is
broken}

\texttt{segment\_polarity\_mc.py} (a multi-cellular extension with
intercellular WG/HH signalling) currently finds zero fixed points under
synchronous update --- it does not reproduce even the single-cell attractors
when embedded in a sheet. Rather than ship a quietly wrong model, it is
marked experimental, excluded from every claim, and listed as open work:
either complete the node set (SMO/CIA/CIR wiring across cell boundaries) or
port the curated id-021 rules to the multi-cell setting.

\subsection{Limitations of scope}

(i) The Turing upshift is measured for one kinetics family and one lattice;
generality is untested. (ii) DCLS is an observational correlation across fly
lines, not a perturbation experiment we performed; causal claims belong to
the original authors' dosage manipulations. (iii) The CF decoder, while
cross-validated over fly lines, is evaluated on one dataset from one
laboratory; a second independent CF-annotated dataset would be the right
next test. (iv) Our $\lambda$ reproduction retains a $\sim 15\%$ gap to the
published central value, attributed to unspecified fitting details in the
original; we chose transparency over tuning.

\subsection{What we would do differently}

Three procedural lessons generalise beyond this project. (i) \emph{Size the
sample to the claim before claiming}: the retracted Study-A surrogate win
existed only at $n = 58$; a power consideration at the outset would have
postponed the claim to $n \approx 170$ and never produced it. (ii)
\emph{Report the majority-class baseline first}: the degenerate basin
classifier of Study B would have been caught in an afternoon by printing one
extra line. (iii) \emph{Fix the window before fitting}: the window
sensitivity of Study C means any single-window $\lambda$ number is
incomplete without its window stated; every $\lambda$ in this report now
carries its window with it.
